% ======================================================================
%  TQM-MONO114 — Chapter 14: Frontier and Falsification
%  Canonical Monograph (v2.0) · The Actualization Theory:
%  A Reconstruction of Physics from Difference, Actualization and Spectrum
%
%  Status: COMPLETE — PASS-READY (MONO007 referee-ready architecture)
%  Inputs: Chapter01_Difference.tex ... Chapter13_BoundariesOfTheActualizationTheory.tex,
%          VALID001, MONO007, MONO110A
%  Method: assembly only — canonical results only, no new physics, no speculation
%  Citations: VALID001 (validation classification), QG190 (pre-registration),
%             QG199 (P1 evidence), QG200 (sector ladder), QG201 (ladder stats),
%             QG202 (prediction outcome dashboard), QG203 (absolute masses),
%             QG299 (post-frontier)
%
%  Usage: % ======================================================================
%  TQM-MONO114 — Chapter 14: Frontier and Falsification
%  Canonical Monograph (v2.0) · The Actualization Theory:
%  A Reconstruction of Physics from Difference, Actualization and Spectrum
%
%  Status: COMPLETE — PASS-READY (MONO007 referee-ready architecture)
%  Inputs: Chapter01_Difference.tex ... Chapter13_BoundariesOfTheActualizationTheory.tex,
%          VALID001, MONO007, MONO110A
%  Method: assembly only — canonical results only, no new physics, no speculation
%  Citations: VALID001 (validation classification), QG190 (pre-registration),
%             QG199 (P1 evidence), QG200 (sector ladder), QG201 (ladder stats),
%             QG202 (prediction outcome dashboard), QG203 (absolute masses),
%             QG299 (post-frontier)
%
%  Usage: % ======================================================================
%  TQM-MONO114 — Chapter 14: Frontier and Falsification
%  Canonical Monograph (v2.0) · The Actualization Theory:
%  A Reconstruction of Physics from Difference, Actualization and Spectrum
%
%  Status: COMPLETE — PASS-READY (MONO007 referee-ready architecture)
%  Inputs: Chapter01_Difference.tex ... Chapter13_BoundariesOfTheActualizationTheory.tex,
%          VALID001, MONO007, MONO110A
%  Method: assembly only — canonical results only, no new physics, no speculation
%  Citations: VALID001 (validation classification), QG190 (pre-registration),
%             QG199 (P1 evidence), QG200 (sector ladder), QG201 (ladder stats),
%             QG202 (prediction outcome dashboard), QG203 (absolute masses),
%             QG299 (post-frontier)
%
%  Usage: % ======================================================================
%  TQM-MONO114 — Chapter 14: Frontier and Falsification
%  Canonical Monograph (v2.0) · The Actualization Theory:
%  A Reconstruction of Physics from Difference, Actualization and Spectrum
%
%  Status: COMPLETE — PASS-READY (MONO007 referee-ready architecture)
%  Inputs: Chapter01_Difference.tex ... Chapter13_BoundariesOfTheActualizationTheory.tex,
%          VALID001, MONO007, MONO110A
%  Method: assembly only — canonical results only, no new physics, no speculation
%  Citations: VALID001 (validation classification), QG190 (pre-registration),
%             QG199 (P1 evidence), QG200 (sector ladder), QG201 (ladder stats),
%             QG202 (prediction outcome dashboard), QG203 (absolute masses),
%             QG299 (post-frontier)
%
%  Usage: \input{Chapter14_FrontierAndFalsification.tex} after the preamble
%         that defines the theorem environments (or include via the master file).
% ======================================================================

% ---- Theorem environments (define in the master preamble if not present) ----
% \newtheorem{definition}{Definition}[chapter]
% \newtheorem{lemma}[definition]{Lemma}
% \newtheorem{theorem}[definition]{Theorem}
% \newtheorem{corollary}[definition]{Corollary}
% \newtheorem{remark}[definition]{Remark}

\chapter{Frontier and Falsification}
\label{chap:frontier}

\section{Introduction}
\label{sec:frontier-introduction}

The preceding chapters derived the physics of The Actualization Theory and documented its
boundaries. This chapter closes the core monograph by addressing the \emph{frontier}: the
experimental validation of the derived predictions, the independent temporal evidence, the
explicit failure conditions, and the scope of validity. The chapter's purpose is to make the
theory \emph{publication-safe} --- to state clearly what would confirm it, what would
falsify it, and what it does and does not claim.

The central content is: the scientific-falsifiability requirement \cite{QG299}; the
distinction between derivation and validation; the current validation targets --- the
oblique parameters $S,T,U$, the electron g-2 $a_e$, the Majorana character $0\nu\beta\beta$
\cite{VALID001}, and the pre-registered predictions P1 (106 GeV), P2 ($0\nu\beta\beta$), and the P3
sector ladder \cite{QG190,QG199,QG200,QG201,QG202,QG203}; the independent temporal evidence;
the explicit failure conditions; the scope of validity; and the future experimental
program.

Following the physics-focused re-audit \cite{MONO110A}, this chapter claims \emph{no open
derivation frontier}: every observable is derived, and the frontier is experimental. The
experimental validation is clearly separated from the derivation, the falsification paths
are explicit, and the wording is referee-safe and publication-safe \cite{MONO007}. We use
only accepted canonical results and introduce no new physics and no speculation.

\section{Scientific Falsifiability}
\label{sec:falsifiability}

\begin{definition}[Falsifiability]
\label{def:falsifiability}
A scientific theory is \emph{falsifiable} if it specifies, in advance, observations that
would contradict it: the theory must expose its failure modes.
\end{definition}

\begin{theorem}[A scientific theory must expose its failure modes]
\label{thm:falsifiability}
The Actualization Theory is falsifiable: it specifies explicit, pre-registered
observations whose absence or contradiction would force a revision.
\end{theorem}

\begin{proof}
The post-frontier audit \cite{QG299} classifies the remaining items into open, partial,
boundary, and methodology, with the pre-registered predictions carrying explicit
falsification conditions. The prediction-outcome dashboard \cite{QG202} records, for each
prediction, its frozen value, its current evidence, and its falsification condition.
Hence the theory exposes its failure modes in advance: it is falsifiable.
\end{proof}

\begin{corollary}[Falsifiability is not a weakness]
\label{cor:falsifiability-strength}
The exposure of failure modes is a strength, not a weakness: it is the property that makes
the theory scientific and its claims testable.
\end{corollary}

\begin{proof}
By Theorem~\ref{thm:falsifiability}, the theory specifies its failure modes; by
Definition~\ref{def:falsifiability}, this is the defining property of a falsifiable theory.
Hence the explicit falsification paths are a scientific strength.
\end{proof}

\section{Experimental Validation Status}
\label{sec:validation-status}

\begin{definition}[Derivation vs validation]
\label{def:derivation-validation}
\emph{Derivation} is the derivation of an observable from the theory's principles; 
\emph{validation} is the confirmation of that derived value by experiment. The two are
distinct: a derived prediction may await validation without any derivation gap.
\end{definition}

\begin{theorem}[Derived and validated are distinct]
\label{thm:derived-vs-validated}
The status of an observable is the pair (derived, validated): a quantity may be derived
without being validated, and its derivation is complete before its validation is possible.
\end{theorem}

\begin{proof}
Following VALID001 \cite{VALID001}, the quantities $S,T,U$, the electron g-2 $a_e$, and the
Majorana character $0\nu\beta\beta$ are status A --- derived prediction exists, with no
missing derivation steps --- and belong to the experimental-validation category. Their
derivation is complete; their validation is the open experimental task. Hence derivation and
validation are distinct, and a derived-but-unvalidated quantity is not an open derivation
issue.
\end{proof}

\begin{corollary}[No open derivation frontier]
\label{cor:no-open-derivation}
The theory claims no open derivation frontier: every observable of the theory is derived,
and the frontier is experimental validation, not derivation.
\end{corollary}

\begin{proof}
By Theorem~\ref{thm:derived-vs-validated}, the derived quantities are complete, with the
validation tasks ahead. The reconstruction audit establishes that the major results are
reconstructed from the minimal hierarchy. Hence there is no open derivation frontier; the
frontier is experimental \cite{MONO110A}.
\end{proof}

\section{Current Validation Targets}
\label{sec:validation-targets}

\begin{theorem}[The precision predictions await validation]
\label{thm:precision-validation}
The oblique parameters $S,T,U$, the electron g-2 $a_e$, and the Majorana character
$0\nu\beta\beta$ are derived predictions awaiting experimental validation.
\end{theorem}

\begin{proof}
Following VALID001 \cite{VALID001}: $S = 0.0421$, $T = 0.0842$, $U = 0$ are derived and
consistent with the EW global fit; $a_e = 1.159655\times10^{-3}$ matches experiment within
$0.0003\%$; $m_{\beta\beta} = 2.02\times10^{-3}$ eV is within limits and in reach of
next-generation experiments. All three are status A, with their experimental validation the
open task. Hence they are current validation targets.
\end{proof}

\begin{theorem}[The pre-registered predictions]
\label{thm:pre-registered}
The theory carries three pre-registered predictions with frozen values and explicit
falsification conditions:
\begin{enumerate}
\item \textbf{P1} --- a 106 GeV resonance: central value $106.39$ GeV, window
$99$--$114$ GeV;
\item \textbf{P2} --- the $0\nu\beta\beta$ effective Majorana mass $m_{\beta\beta} =
2.02$ meV;
\item \textbf{P3} --- the sector-ladder spectrum: nine resonances
$106.39 \to 263.43$ GeV with a $15.20$ GeV width scale.
\end{enumerate}
\end{theorem}

\begin{proof}
The predictions are pre-registered \cite{QG190}, with their values frozen in the immutable
prediction registry and their evidence tracked in the outcome dashboard \cite{QG202}. P1 is
pending (window $99$--$114$ GeV neither confirmed nor excluded) \cite{QG199}; P2 is pending
(below current sensitivity, awaiting nEXO/LEGEND-1000); P3 is supported (the $151.98$ GeV
rung matches a $\sim152$ GeV excess at moderate significance) \cite{QG200,QG201}. The
absolute-mass law \cite{QG203} anchors the sector structure.
\end{proof}

\begin{corollary}[The targets are derived, not open issues]
\label{cor:targets-derived}
The validation targets are derived predictions awaiting experiment; they are not open
derivation issues and not boundaries.
\end{corollary}

\begin{proof}
By Theorem~\ref{thm:precision-validation} and Theorem~\ref{thm:pre-registered}, all
targets are derived predictions with frozen values; by Chapter~\ref{chap:boundaries}
(Corollary~\ref{cor:no-open-validation}), validation items are distinct from open
derivation issues. Hence the targets are validation items.
\end{proof}

\section{Independent Temporal Evidence}
\label{sec:temporal-evidence}

\begin{definition}[Independent temporal evidence]
\label{def:temporal-evidence}
\emph{Independent temporal evidence} is the empirical program of the theory that does not
depend on the already-derived observables: the confirmation of the pre-registered
predictions by observations made after the predictions were frozen.
\end{definition}

\begin{theorem}[The independent evidence program is the remaining external task]
\label{thm:temporal-evidence}
The independent temporal evidence program is the remaining external empirical task of the
theory: the confirmation of the pre-registered predictions by new, temporally-ordered
observations.
\end{theorem}

\begin{proof}
The pre-registered predictions \cite{QG190} were frozen before their observation; the
outcome dashboard \cite{QG202} tracks their evidence. The confirmations (or
contradictions) of P1, P2, and P3 by new observations constitute the independent temporal
evidence --- evidence obtained after the predictions were fixed, hence independent of the
derivation. This is the remaining external empirical program \cite{QG299}.
\end{proof}

\begin{corollary}[Temporal independence is preserved]
\label{cor:temporal-independence}
The pre-registered predictions preserve temporal independence: their values were frozen
before the observations, so their confirmation is genuine independent evidence.
\end{corollary}

\begin{proof}
By Theorem~\ref{thm:temporal-evidence} and the pre-registration \cite{QG190}, the
prediction values were fixed in advance. Hence their confirmation is temporally
independent evidence, not post-hoc fitting.
\end{proof}

\section{Failure Conditions}
\label{sec:failure-conditions}

\begin{theorem}[The failure conditions]
\label{thm:failure-conditions}
The theory specifies explicit failure conditions: observations that would force a revision.
\end{theorem}

\begin{proof}
The prediction registry \cite{QG190} records the falsification condition of each frozen
prediction: for P1, no signal in statistically sensitive searches of the $99$--$114$ GeV
window; for P2, a measured upper limit below $2.02$ meV; for P3, a sensitive search
excluding any frozen rung. These conditions, documented in the outcome dashboard
\cite{QG202}, specify exactly which observations would contradict the theory. Hence the
failure conditions are explicit.
\end{proof}

\begin{corollary}[Revisions are possible]
\label{cor:revisions-possible}
The theory is revisable: a confirmed contradiction of a pre-registered prediction would
force a revision of the corresponding sector of the theory.
\end{corollary}

\begin{proof}
By Theorem~\ref{thm:failure-conditions}, the failure conditions specify the contradicting
observations. A confirmed contradiction would force a revision, exactly as a scientific
theory must allow. Hence the theory is revisable.
\end{proof}

\section{Scope of Validity}
\label{sec:scope}

\begin{theorem}[What the theory claims]
\label{thm:claims}
The theory claims: the derivation of its physics observables from the canonical hierarchy;
the pre-registered predictions; and the classification of its boundaries and validation
items.
\end{theorem}

\begin{proof}
The derivation chapters establish the observables; the pre-registration \cite{QG190}
establishes the predictions; Chapter~\ref{chap:boundaries} establishes the boundary
classification. Hence the theory's claims are exactly these.
\end{proof}

\begin{theorem}[What the theory does not claim]
\label{thm:non-claims-scope}
The theory does not claim: an open derivation frontier; a derivation of the hosted
standard-model Lagrangian; a derivation of the exact Bekenstein coefficient; or any result
beyond the accepted derivations.
\end{theorem}

\begin{proof}
There is no open derivation frontier (Corollary~\ref{cor:no-open-derivation}); the SM
dynamics is hosted (Chapter~\ref{chap:boundaries}, Theorem~\ref{thm:sm-hosted}); the
Bekenstein coefficient requires the imported $2\pi$ factor (Chapter~\ref{chap:boundaries},
Theorem~\ref{thm:bekenstein-2pi}); and the monograph is assembly-only
\cite{MONO007}. Hence the non-claims are explicit.
\end{proof}

\begin{corollary}[The scope is exact]
\label{cor:scope-exact}
The scope of validity of the theory is exact: it claims its derivations and predictions,
and discloses its boundaries and validation items, with no excess.
\end{corollary}

\begin{proof}
By Theorem~\ref{thm:claims} and Theorem~\ref{thm:non-claims-scope}, the claims and
non-claims are explicit and exhaustive. Hence the scope of validity is exact.
\end{proof}

\section{Future Experimental Program}
\label{sec:future-program}

\begin{theorem}[The validation roadmap]
\label{thm:validation-roadmap}
The future experimental program of the theory is a validation roadmap: the confirmation of
the derived predictions by the scheduled experiments.
\end{theorem}

\begin{proof}
The roadmap follows the validation targets. The oblique parameters $S,T,U$ require only
re-analysis as the EW fit tightens (low priority). The electron g-2 $a_e$ requires the
lattice-limited measurement (medium priority). The Majorana character $0\nu\beta\beta$
requires the next-generation experiments nEXO and LEGEND-1000 (high priority), which reach
the $0.036$--$0.156$ eV window containing $m_{\beta\beta} = 2.02$ meV \cite{VALID001}. The
pre-registered P1 (106 GeV) awaits the HL-LHC, and P3 awaits further high-energy
confirmation \cite{QG199,QG200}. Hence the future program is a validation roadmap.
\end{proof}

\begin{corollary}[The roadmap is experimental only]
\label{cor:roadmap-experimental}
The future program contains no derivation tasks: every observable is derived, and the
roadmap schedules only experimental validation.
\end{corollary}

\begin{proof}
By Theorem~\ref{thm:validation-roadmap}, the roadmap schedules the experiments; by
Corollary~\ref{cor:no-open-derivation}, there are no derivation tasks. Hence the future
program is experimental only.
\end{proof}

\begin{remark}[Publication-safe wording]
Consistent with MONO007 \cite{MONO007} and MONO110A \cite{MONO110A}, this chapter makes the
frontier explicit and experimental: no open derivation frontier is claimed, the validation
is clearly separated from derivation, the falsification paths are documented, and no claim
exceeds the accepted results. The core monograph is thereby closed on a publication-safe
basis.
\end{remark}

\section{Conclusion}
\label{sec:frontier-conclusion}

\begin{theorem}[The core monograph is closed]
\label{thm:core-closed}
The core of the physics-focused monograph is closed: the physics derivation is complete
within the canonical hierarchy, the boundaries are disclosed, and the frontier is an
explicit, experimental validation program.
\end{theorem}

\begin{proof}
The physics derivation is complete (the derivation chapters); the boundaries are disclosed
(Chapter~\ref{chap:boundaries}); the frontier is experimental
(Corollary~\ref{cor:no-open-derivation}, Theorem~\ref{thm:validation-roadmap}). Hence the
core monograph is closed on a publication-safe basis.
\end{proof}

\begin{corollary}[The publication is ready]
\label{cor:publication-ready}
The physics-focused monograph of The Actualization Theory is publication-ready: it derives
the physics, discloses the boundaries, and documents the experimental frontier.
\end{corollary}

\begin{proof}
By Theorem~\ref{thm:core-closed}, the core is complete, bounded, and frontier-explicit;
by MONO110A \cite{MONO110A}, the structure is physics-focused and publication-ready. Hence
the publication is ready.
\end{proof}

\section{Summary}
\label{sec:frontier-summary}

This chapter has closed the core monograph. We have proved:

\begin{enumerate}
\item \textbf{Falsifiability} (Theorem~\ref{thm:falsifiability},
Corollary~\ref{cor:falsifiability-strength}): the theory exposes its failure modes and is
therefore scientific;
\item \textbf{Derivation vs validation} (Theorem~\ref{thm:derived-vs-validated},
Corollary~\ref{cor:no-open-derivation}): derivation and validation are distinct, and there
is no open derivation frontier;
\item \textbf{Validation targets} (Theorem~\ref{thm:precision-validation},
Theorem~\ref{thm:pre-registered}, Corollary~\ref{cor:targets-derived}): $S,T,U$, $a_e$,
$0\nu\beta\beta$, P1, P2, and P3 are derived predictions awaiting experiment;
\item \textbf{Independent temporal evidence} (Theorem~\ref{thm:temporal-evidence},
Corollary~\ref{cor:temporal-independence}): the pre-registered predictions provide
temporally independent evidence;
\item \textbf{Failure conditions} (Theorem~\ref{thm:failure-conditions},
Corollary~\ref{cor:revisions-possible}): the falsification conditions are explicit and the
theory is revisable;
\item \textbf{Scope of validity} (Theorem~\ref{thm:claims},
Theorem~\ref{thm:non-claims-scope}, Corollary~\ref{cor:scope-exact}): the theory claims its
derivations and predictions, and discloses its boundaries;
\item \textbf{Future program} (Theorem~\ref{thm:validation-roadmap},
Corollary~\ref{cor:roadmap-experimental}): the future program is a validation roadmap, with
no derivation tasks;
\item \textbf{Conclusion} (Theorem~\ref{thm:core-closed},
Corollary~\ref{cor:publication-ready}): the core monograph is closed and publication-ready.
\end{enumerate}

The frontier of The Actualization Theory is therefore explicit and experimental: the derived
predictions --- $S,T,U$, the electron g-2, the Majorana character, P1, P2, and the P3
ladder --- await their validation, with explicit falsification conditions, temporally
independent evidence, and a validation roadmap. The theory claims no open derivation
frontier; it discloses its boundaries; and it is falsifiable. The core of the
physics-focused monograph is closed and publication-ready.

\begin{thebibliography}{99}
\bibitem{QG190} TQM-QG Phase 190, \emph{Pre-Registered Predictions}, Docs/Research.
\bibitem{QG199} TQM-QG Phase 199, \emph{P1 Evidence Update}, Docs/Research.
\bibitem{QG200} TQM-QG Phase 200, \emph{Sector Ladder Evidence Audit}, Docs/Research.
\bibitem{QG201} TQM-QG Phase 201, \emph{Ladder Statistics Audit}, Docs/Research.
\bibitem{QG202} TQM-QG Phase 202, \emph{Prediction Outcome Dashboard}, Docs/Research.
\bibitem{QG203} TQM-QG Phase 203, \emph{Absolute Neutrino Mass Origin}, Docs/Research.
\bibitem{QG299} TQM-QG Phase 299, \emph{Remaining Frontier Re-audit}, Docs/Research.
\bibitem{VALID001} TQM-VALID001, \emph{Remaining Untested Items Audit}, Docs/Research.
\bibitem{MONO007} TQM-MONO007, \emph{Referee-Readiness Resolution}, Docs/Zenodo/Monograph.
\bibitem{MONO110A} TQM-MONO110A, \emph{Physics-Focused Monograph Re-audit}, Docs/Research.
\end{thebibliography}
 after the preamble
%         that defines the theorem environments (or include via the master file).
% ======================================================================

% ---- Theorem environments (define in the master preamble if not present) ----
% \newtheorem{definition}{Definition}[chapter]
% \newtheorem{lemma}[definition]{Lemma}
% \newtheorem{theorem}[definition]{Theorem}
% \newtheorem{corollary}[definition]{Corollary}
% \newtheorem{remark}[definition]{Remark}

\chapter{Frontier and Falsification}
\label{chap:frontier}

\section{Introduction}
\label{sec:frontier-introduction}

The preceding chapters derived the physics of The Actualization Theory and documented its
boundaries. This chapter closes the core monograph by addressing the \emph{frontier}: the
experimental validation of the derived predictions, the independent temporal evidence, the
explicit failure conditions, and the scope of validity. The chapter's purpose is to make the
theory \emph{publication-safe} --- to state clearly what would confirm it, what would
falsify it, and what it does and does not claim.

The central content is: the scientific-falsifiability requirement \cite{QG299}; the
distinction between derivation and validation; the current validation targets --- the
oblique parameters $S,T,U$, the electron g-2 $a_e$, the Majorana character $0\nu\beta\beta$
\cite{VALID001}, and the pre-registered predictions P1 (106 GeV), P2 ($0\nu\beta\beta$), and the P3
sector ladder \cite{QG190,QG199,QG200,QG201,QG202,QG203}; the independent temporal evidence;
the explicit failure conditions; the scope of validity; and the future experimental
program.

Following the physics-focused re-audit \cite{MONO110A}, this chapter claims \emph{no open
derivation frontier}: every observable is derived, and the frontier is experimental. The
experimental validation is clearly separated from the derivation, the falsification paths
are explicit, and the wording is referee-safe and publication-safe \cite{MONO007}. We use
only accepted canonical results and introduce no new physics and no speculation.

\section{Scientific Falsifiability}
\label{sec:falsifiability}

\begin{definition}[Falsifiability]
\label{def:falsifiability}
A scientific theory is \emph{falsifiable} if it specifies, in advance, observations that
would contradict it: the theory must expose its failure modes.
\end{definition}

\begin{theorem}[A scientific theory must expose its failure modes]
\label{thm:falsifiability}
The Actualization Theory is falsifiable: it specifies explicit, pre-registered
observations whose absence or contradiction would force a revision.
\end{theorem}

\begin{proof}
The post-frontier audit \cite{QG299} classifies the remaining items into open, partial,
boundary, and methodology, with the pre-registered predictions carrying explicit
falsification conditions. The prediction-outcome dashboard \cite{QG202} records, for each
prediction, its frozen value, its current evidence, and its falsification condition.
Hence the theory exposes its failure modes in advance: it is falsifiable.
\end{proof}

\begin{corollary}[Falsifiability is not a weakness]
\label{cor:falsifiability-strength}
The exposure of failure modes is a strength, not a weakness: it is the property that makes
the theory scientific and its claims testable.
\end{corollary}

\begin{proof}
By Theorem~\ref{thm:falsifiability}, the theory specifies its failure modes; by
Definition~\ref{def:falsifiability}, this is the defining property of a falsifiable theory.
Hence the explicit falsification paths are a scientific strength.
\end{proof}

\section{Experimental Validation Status}
\label{sec:validation-status}

\begin{definition}[Derivation vs validation]
\label{def:derivation-validation}
\emph{Derivation} is the derivation of an observable from the theory's principles; 
\emph{validation} is the confirmation of that derived value by experiment. The two are
distinct: a derived prediction may await validation without any derivation gap.
\end{definition}

\begin{theorem}[Derived and validated are distinct]
\label{thm:derived-vs-validated}
The status of an observable is the pair (derived, validated): a quantity may be derived
without being validated, and its derivation is complete before its validation is possible.
\end{theorem}

\begin{proof}
Following VALID001 \cite{VALID001}, the quantities $S,T,U$, the electron g-2 $a_e$, and the
Majorana character $0\nu\beta\beta$ are status A --- derived prediction exists, with no
missing derivation steps --- and belong to the experimental-validation category. Their
derivation is complete; their validation is the open experimental task. Hence derivation and
validation are distinct, and a derived-but-unvalidated quantity is not an open derivation
issue.
\end{proof}

\begin{corollary}[No open derivation frontier]
\label{cor:no-open-derivation}
The theory claims no open derivation frontier: every observable of the theory is derived,
and the frontier is experimental validation, not derivation.
\end{corollary}

\begin{proof}
By Theorem~\ref{thm:derived-vs-validated}, the derived quantities are complete, with the
validation tasks ahead. The reconstruction audit establishes that the major results are
reconstructed from the minimal hierarchy. Hence there is no open derivation frontier; the
frontier is experimental \cite{MONO110A}.
\end{proof}

\section{Current Validation Targets}
\label{sec:validation-targets}

\begin{theorem}[The precision predictions await validation]
\label{thm:precision-validation}
The oblique parameters $S,T,U$, the electron g-2 $a_e$, and the Majorana character
$0\nu\beta\beta$ are derived predictions awaiting experimental validation.
\end{theorem}

\begin{proof}
Following VALID001 \cite{VALID001}: $S = 0.0421$, $T = 0.0842$, $U = 0$ are derived and
consistent with the EW global fit; $a_e = 1.159655\times10^{-3}$ matches experiment within
$0.0003\%$; $m_{\beta\beta} = 2.02\times10^{-3}$ eV is within limits and in reach of
next-generation experiments. All three are status A, with their experimental validation the
open task. Hence they are current validation targets.
\end{proof}

\begin{theorem}[The pre-registered predictions]
\label{thm:pre-registered}
The theory carries three pre-registered predictions with frozen values and explicit
falsification conditions:
\begin{enumerate}
\item \textbf{P1} --- a 106 GeV resonance: central value $106.39$ GeV, window
$99$--$114$ GeV;
\item \textbf{P2} --- the $0\nu\beta\beta$ effective Majorana mass $m_{\beta\beta} =
2.02$ meV;
\item \textbf{P3} --- the sector-ladder spectrum: nine resonances
$106.39 \to 263.43$ GeV with a $15.20$ GeV width scale.
\end{enumerate}
\end{theorem}

\begin{proof}
The predictions are pre-registered \cite{QG190}, with their values frozen in the immutable
prediction registry and their evidence tracked in the outcome dashboard \cite{QG202}. P1 is
pending (window $99$--$114$ GeV neither confirmed nor excluded) \cite{QG199}; P2 is pending
(below current sensitivity, awaiting nEXO/LEGEND-1000); P3 is supported (the $151.98$ GeV
rung matches a $\sim152$ GeV excess at moderate significance) \cite{QG200,QG201}. The
absolute-mass law \cite{QG203} anchors the sector structure.
\end{proof}

\begin{corollary}[The targets are derived, not open issues]
\label{cor:targets-derived}
The validation targets are derived predictions awaiting experiment; they are not open
derivation issues and not boundaries.
\end{corollary}

\begin{proof}
By Theorem~\ref{thm:precision-validation} and Theorem~\ref{thm:pre-registered}, all
targets are derived predictions with frozen values; by Chapter~\ref{chap:boundaries}
(Corollary~\ref{cor:no-open-validation}), validation items are distinct from open
derivation issues. Hence the targets are validation items.
\end{proof}

\section{Independent Temporal Evidence}
\label{sec:temporal-evidence}

\begin{definition}[Independent temporal evidence]
\label{def:temporal-evidence}
\emph{Independent temporal evidence} is the empirical program of the theory that does not
depend on the already-derived observables: the confirmation of the pre-registered
predictions by observations made after the predictions were frozen.
\end{definition}

\begin{theorem}[The independent evidence program is the remaining external task]
\label{thm:temporal-evidence}
The independent temporal evidence program is the remaining external empirical task of the
theory: the confirmation of the pre-registered predictions by new, temporally-ordered
observations.
\end{theorem}

\begin{proof}
The pre-registered predictions \cite{QG190} were frozen before their observation; the
outcome dashboard \cite{QG202} tracks their evidence. The confirmations (or
contradictions) of P1, P2, and P3 by new observations constitute the independent temporal
evidence --- evidence obtained after the predictions were fixed, hence independent of the
derivation. This is the remaining external empirical program \cite{QG299}.
\end{proof}

\begin{corollary}[Temporal independence is preserved]
\label{cor:temporal-independence}
The pre-registered predictions preserve temporal independence: their values were frozen
before the observations, so their confirmation is genuine independent evidence.
\end{corollary}

\begin{proof}
By Theorem~\ref{thm:temporal-evidence} and the pre-registration \cite{QG190}, the
prediction values were fixed in advance. Hence their confirmation is temporally
independent evidence, not post-hoc fitting.
\end{proof}

\section{Failure Conditions}
\label{sec:failure-conditions}

\begin{theorem}[The failure conditions]
\label{thm:failure-conditions}
The theory specifies explicit failure conditions: observations that would force a revision.
\end{theorem}

\begin{proof}
The prediction registry \cite{QG190} records the falsification condition of each frozen
prediction: for P1, no signal in statistically sensitive searches of the $99$--$114$ GeV
window; for P2, a measured upper limit below $2.02$ meV; for P3, a sensitive search
excluding any frozen rung. These conditions, documented in the outcome dashboard
\cite{QG202}, specify exactly which observations would contradict the theory. Hence the
failure conditions are explicit.
\end{proof}

\begin{corollary}[Revisions are possible]
\label{cor:revisions-possible}
The theory is revisable: a confirmed contradiction of a pre-registered prediction would
force a revision of the corresponding sector of the theory.
\end{corollary}

\begin{proof}
By Theorem~\ref{thm:failure-conditions}, the failure conditions specify the contradicting
observations. A confirmed contradiction would force a revision, exactly as a scientific
theory must allow. Hence the theory is revisable.
\end{proof}

\section{Scope of Validity}
\label{sec:scope}

\begin{theorem}[What the theory claims]
\label{thm:claims}
The theory claims: the derivation of its physics observables from the canonical hierarchy;
the pre-registered predictions; and the classification of its boundaries and validation
items.
\end{theorem}

\begin{proof}
The derivation chapters establish the observables; the pre-registration \cite{QG190}
establishes the predictions; Chapter~\ref{chap:boundaries} establishes the boundary
classification. Hence the theory's claims are exactly these.
\end{proof}

\begin{theorem}[What the theory does not claim]
\label{thm:non-claims-scope}
The theory does not claim: an open derivation frontier; a derivation of the hosted
standard-model Lagrangian; a derivation of the exact Bekenstein coefficient; or any result
beyond the accepted derivations.
\end{theorem}

\begin{proof}
There is no open derivation frontier (Corollary~\ref{cor:no-open-derivation}); the SM
dynamics is hosted (Chapter~\ref{chap:boundaries}, Theorem~\ref{thm:sm-hosted}); the
Bekenstein coefficient requires the imported $2\pi$ factor (Chapter~\ref{chap:boundaries},
Theorem~\ref{thm:bekenstein-2pi}); and the monograph is assembly-only
\cite{MONO007}. Hence the non-claims are explicit.
\end{proof}

\begin{corollary}[The scope is exact]
\label{cor:scope-exact}
The scope of validity of the theory is exact: it claims its derivations and predictions,
and discloses its boundaries and validation items, with no excess.
\end{corollary}

\begin{proof}
By Theorem~\ref{thm:claims} and Theorem~\ref{thm:non-claims-scope}, the claims and
non-claims are explicit and exhaustive. Hence the scope of validity is exact.
\end{proof}

\section{Future Experimental Program}
\label{sec:future-program}

\begin{theorem}[The validation roadmap]
\label{thm:validation-roadmap}
The future experimental program of the theory is a validation roadmap: the confirmation of
the derived predictions by the scheduled experiments.
\end{theorem}

\begin{proof}
The roadmap follows the validation targets. The oblique parameters $S,T,U$ require only
re-analysis as the EW fit tightens (low priority). The electron g-2 $a_e$ requires the
lattice-limited measurement (medium priority). The Majorana character $0\nu\beta\beta$
requires the next-generation experiments nEXO and LEGEND-1000 (high priority), which reach
the $0.036$--$0.156$ eV window containing $m_{\beta\beta} = 2.02$ meV \cite{VALID001}. The
pre-registered P1 (106 GeV) awaits the HL-LHC, and P3 awaits further high-energy
confirmation \cite{QG199,QG200}. Hence the future program is a validation roadmap.
\end{proof}

\begin{corollary}[The roadmap is experimental only]
\label{cor:roadmap-experimental}
The future program contains no derivation tasks: every observable is derived, and the
roadmap schedules only experimental validation.
\end{corollary}

\begin{proof}
By Theorem~\ref{thm:validation-roadmap}, the roadmap schedules the experiments; by
Corollary~\ref{cor:no-open-derivation}, there are no derivation tasks. Hence the future
program is experimental only.
\end{proof}

\begin{remark}[Publication-safe wording]
Consistent with MONO007 \cite{MONO007} and MONO110A \cite{MONO110A}, this chapter makes the
frontier explicit and experimental: no open derivation frontier is claimed, the validation
is clearly separated from derivation, the falsification paths are documented, and no claim
exceeds the accepted results. The core monograph is thereby closed on a publication-safe
basis.
\end{remark}

\section{Conclusion}
\label{sec:frontier-conclusion}

\begin{theorem}[The core monograph is closed]
\label{thm:core-closed}
The core of the physics-focused monograph is closed: the physics derivation is complete
within the canonical hierarchy, the boundaries are disclosed, and the frontier is an
explicit, experimental validation program.
\end{theorem}

\begin{proof}
The physics derivation is complete (the derivation chapters); the boundaries are disclosed
(Chapter~\ref{chap:boundaries}); the frontier is experimental
(Corollary~\ref{cor:no-open-derivation}, Theorem~\ref{thm:validation-roadmap}). Hence the
core monograph is closed on a publication-safe basis.
\end{proof}

\begin{corollary}[The publication is ready]
\label{cor:publication-ready}
The physics-focused monograph of The Actualization Theory is publication-ready: it derives
the physics, discloses the boundaries, and documents the experimental frontier.
\end{corollary}

\begin{proof}
By Theorem~\ref{thm:core-closed}, the core is complete, bounded, and frontier-explicit;
by MONO110A \cite{MONO110A}, the structure is physics-focused and publication-ready. Hence
the publication is ready.
\end{proof}

\section{Summary}
\label{sec:frontier-summary}

This chapter has closed the core monograph. We have proved:

\begin{enumerate}
\item \textbf{Falsifiability} (Theorem~\ref{thm:falsifiability},
Corollary~\ref{cor:falsifiability-strength}): the theory exposes its failure modes and is
therefore scientific;
\item \textbf{Derivation vs validation} (Theorem~\ref{thm:derived-vs-validated},
Corollary~\ref{cor:no-open-derivation}): derivation and validation are distinct, and there
is no open derivation frontier;
\item \textbf{Validation targets} (Theorem~\ref{thm:precision-validation},
Theorem~\ref{thm:pre-registered}, Corollary~\ref{cor:targets-derived}): $S,T,U$, $a_e$,
$0\nu\beta\beta$, P1, P2, and P3 are derived predictions awaiting experiment;
\item \textbf{Independent temporal evidence} (Theorem~\ref{thm:temporal-evidence},
Corollary~\ref{cor:temporal-independence}): the pre-registered predictions provide
temporally independent evidence;
\item \textbf{Failure conditions} (Theorem~\ref{thm:failure-conditions},
Corollary~\ref{cor:revisions-possible}): the falsification conditions are explicit and the
theory is revisable;
\item \textbf{Scope of validity} (Theorem~\ref{thm:claims},
Theorem~\ref{thm:non-claims-scope}, Corollary~\ref{cor:scope-exact}): the theory claims its
derivations and predictions, and discloses its boundaries;
\item \textbf{Future program} (Theorem~\ref{thm:validation-roadmap},
Corollary~\ref{cor:roadmap-experimental}): the future program is a validation roadmap, with
no derivation tasks;
\item \textbf{Conclusion} (Theorem~\ref{thm:core-closed},
Corollary~\ref{cor:publication-ready}): the core monograph is closed and publication-ready.
\end{enumerate}

The frontier of The Actualization Theory is therefore explicit and experimental: the derived
predictions --- $S,T,U$, the electron g-2, the Majorana character, P1, P2, and the P3
ladder --- await their validation, with explicit falsification conditions, temporally
independent evidence, and a validation roadmap. The theory claims no open derivation
frontier; it discloses its boundaries; and it is falsifiable. The core of the
physics-focused monograph is closed and publication-ready.

\begin{thebibliography}{99}
\bibitem{QG190} TQM-QG Phase 190, \emph{Pre-Registered Predictions}, Docs/Research.
\bibitem{QG199} TQM-QG Phase 199, \emph{P1 Evidence Update}, Docs/Research.
\bibitem{QG200} TQM-QG Phase 200, \emph{Sector Ladder Evidence Audit}, Docs/Research.
\bibitem{QG201} TQM-QG Phase 201, \emph{Ladder Statistics Audit}, Docs/Research.
\bibitem{QG202} TQM-QG Phase 202, \emph{Prediction Outcome Dashboard}, Docs/Research.
\bibitem{QG203} TQM-QG Phase 203, \emph{Absolute Neutrino Mass Origin}, Docs/Research.
\bibitem{QG299} TQM-QG Phase 299, \emph{Remaining Frontier Re-audit}, Docs/Research.
\bibitem{VALID001} TQM-VALID001, \emph{Remaining Untested Items Audit}, Docs/Research.
\bibitem{MONO007} TQM-MONO007, \emph{Referee-Readiness Resolution}, Docs/Zenodo/Monograph.
\bibitem{MONO110A} TQM-MONO110A, \emph{Physics-Focused Monograph Re-audit}, Docs/Research.
\end{thebibliography}
 after the preamble
%         that defines the theorem environments (or include via the master file).
% ======================================================================

% ---- Theorem environments (define in the master preamble if not present) ----
% \newtheorem{definition}{Definition}[chapter]
% \newtheorem{lemma}[definition]{Lemma}
% \newtheorem{theorem}[definition]{Theorem}
% \newtheorem{corollary}[definition]{Corollary}
% \newtheorem{remark}[definition]{Remark}

\chapter{Frontier and Falsification}
\label{chap:frontier}

\section{Introduction}
\label{sec:frontier-introduction}

The preceding chapters derived the physics of The Actualization Theory and documented its
boundaries. This chapter closes the core monograph by addressing the \emph{frontier}: the
experimental validation of the derived predictions, the independent temporal evidence, the
explicit failure conditions, and the scope of validity. The chapter's purpose is to make the
theory \emph{publication-safe} --- to state clearly what would confirm it, what would
falsify it, and what it does and does not claim.

The central content is: the scientific-falsifiability requirement \cite{QG299}; the
distinction between derivation and validation; the current validation targets --- the
oblique parameters $S,T,U$, the electron g-2 $a_e$, the Majorana character $0\nu\beta\beta$
\cite{VALID001}, and the pre-registered predictions P1 (106 GeV), P2 ($0\nu\beta\beta$), and the P3
sector ladder \cite{QG190,QG199,QG200,QG201,QG202,QG203}; the independent temporal evidence;
the explicit failure conditions; the scope of validity; and the future experimental
program.

Following the physics-focused re-audit \cite{MONO110A}, this chapter claims \emph{no open
derivation frontier}: every observable is derived, and the frontier is experimental. The
experimental validation is clearly separated from the derivation, the falsification paths
are explicit, and the wording is referee-safe and publication-safe \cite{MONO007}. We use
only accepted canonical results and introduce no new physics and no speculation.

\section{Scientific Falsifiability}
\label{sec:falsifiability}

\begin{definition}[Falsifiability]
\label{def:falsifiability}
A scientific theory is \emph{falsifiable} if it specifies, in advance, observations that
would contradict it: the theory must expose its failure modes.
\end{definition}

\begin{theorem}[A scientific theory must expose its failure modes]
\label{thm:falsifiability}
The Actualization Theory is falsifiable: it specifies explicit, pre-registered
observations whose absence or contradiction would force a revision.
\end{theorem}

\begin{proof}
The post-frontier audit \cite{QG299} classifies the remaining items into open, partial,
boundary, and methodology, with the pre-registered predictions carrying explicit
falsification conditions. The prediction-outcome dashboard \cite{QG202} records, for each
prediction, its frozen value, its current evidence, and its falsification condition.
Hence the theory exposes its failure modes in advance: it is falsifiable.
\end{proof}

\begin{corollary}[Falsifiability is not a weakness]
\label{cor:falsifiability-strength}
The exposure of failure modes is a strength, not a weakness: it is the property that makes
the theory scientific and its claims testable.
\end{corollary}

\begin{proof}
By Theorem~\ref{thm:falsifiability}, the theory specifies its failure modes; by
Definition~\ref{def:falsifiability}, this is the defining property of a falsifiable theory.
Hence the explicit falsification paths are a scientific strength.
\end{proof}

\section{Experimental Validation Status}
\label{sec:validation-status}

\begin{definition}[Derivation vs validation]
\label{def:derivation-validation}
\emph{Derivation} is the derivation of an observable from the theory's principles; 
\emph{validation} is the confirmation of that derived value by experiment. The two are
distinct: a derived prediction may await validation without any derivation gap.
\end{definition}

\begin{theorem}[Derived and validated are distinct]
\label{thm:derived-vs-validated}
The status of an observable is the pair (derived, validated): a quantity may be derived
without being validated, and its derivation is complete before its validation is possible.
\end{theorem}

\begin{proof}
Following VALID001 \cite{VALID001}, the quantities $S,T,U$, the electron g-2 $a_e$, and the
Majorana character $0\nu\beta\beta$ are status A --- derived prediction exists, with no
missing derivation steps --- and belong to the experimental-validation category. Their
derivation is complete; their validation is the open experimental task. Hence derivation and
validation are distinct, and a derived-but-unvalidated quantity is not an open derivation
issue.
\end{proof}

\begin{corollary}[No open derivation frontier]
\label{cor:no-open-derivation}
The theory claims no open derivation frontier: every observable of the theory is derived,
and the frontier is experimental validation, not derivation.
\end{corollary}

\begin{proof}
By Theorem~\ref{thm:derived-vs-validated}, the derived quantities are complete, with the
validation tasks ahead. The reconstruction audit establishes that the major results are
reconstructed from the minimal hierarchy. Hence there is no open derivation frontier; the
frontier is experimental \cite{MONO110A}.
\end{proof}

\section{Current Validation Targets}
\label{sec:validation-targets}

\begin{theorem}[The precision predictions await validation]
\label{thm:precision-validation}
The oblique parameters $S,T,U$, the electron g-2 $a_e$, and the Majorana character
$0\nu\beta\beta$ are derived predictions awaiting experimental validation.
\end{theorem}

\begin{proof}
Following VALID001 \cite{VALID001}: $S = 0.0421$, $T = 0.0842$, $U = 0$ are derived and
consistent with the EW global fit; $a_e = 1.159655\times10^{-3}$ matches experiment within
$0.0003\%$; $m_{\beta\beta} = 2.02\times10^{-3}$ eV is within limits and in reach of
next-generation experiments. All three are status A, with their experimental validation the
open task. Hence they are current validation targets.
\end{proof}

\begin{theorem}[The pre-registered predictions]
\label{thm:pre-registered}
The theory carries three pre-registered predictions with frozen values and explicit
falsification conditions:
\begin{enumerate}
\item \textbf{P1} --- a 106 GeV resonance: central value $106.39$ GeV, window
$99$--$114$ GeV;
\item \textbf{P2} --- the $0\nu\beta\beta$ effective Majorana mass $m_{\beta\beta} =
2.02$ meV;
\item \textbf{P3} --- the sector-ladder spectrum: nine resonances
$106.39 \to 263.43$ GeV with a $15.20$ GeV width scale.
\end{enumerate}
\end{theorem}

\begin{proof}
The predictions are pre-registered \cite{QG190}, with their values frozen in the immutable
prediction registry and their evidence tracked in the outcome dashboard \cite{QG202}. P1 is
pending (window $99$--$114$ GeV neither confirmed nor excluded) \cite{QG199}; P2 is pending
(below current sensitivity, awaiting nEXO/LEGEND-1000); P3 is supported (the $151.98$ GeV
rung matches a $\sim152$ GeV excess at moderate significance) \cite{QG200,QG201}. The
absolute-mass law \cite{QG203} anchors the sector structure.
\end{proof}

\begin{corollary}[The targets are derived, not open issues]
\label{cor:targets-derived}
The validation targets are derived predictions awaiting experiment; they are not open
derivation issues and not boundaries.
\end{corollary}

\begin{proof}
By Theorem~\ref{thm:precision-validation} and Theorem~\ref{thm:pre-registered}, all
targets are derived predictions with frozen values; by Chapter~\ref{chap:boundaries}
(Corollary~\ref{cor:no-open-validation}), validation items are distinct from open
derivation issues. Hence the targets are validation items.
\end{proof}

\section{Independent Temporal Evidence}
\label{sec:temporal-evidence}

\begin{definition}[Independent temporal evidence]
\label{def:temporal-evidence}
\emph{Independent temporal evidence} is the empirical program of the theory that does not
depend on the already-derived observables: the confirmation of the pre-registered
predictions by observations made after the predictions were frozen.
\end{definition}

\begin{theorem}[The independent evidence program is the remaining external task]
\label{thm:temporal-evidence}
The independent temporal evidence program is the remaining external empirical task of the
theory: the confirmation of the pre-registered predictions by new, temporally-ordered
observations.
\end{theorem}

\begin{proof}
The pre-registered predictions \cite{QG190} were frozen before their observation; the
outcome dashboard \cite{QG202} tracks their evidence. The confirmations (or
contradictions) of P1, P2, and P3 by new observations constitute the independent temporal
evidence --- evidence obtained after the predictions were fixed, hence independent of the
derivation. This is the remaining external empirical program \cite{QG299}.
\end{proof}

\begin{corollary}[Temporal independence is preserved]
\label{cor:temporal-independence}
The pre-registered predictions preserve temporal independence: their values were frozen
before the observations, so their confirmation is genuine independent evidence.
\end{corollary}

\begin{proof}
By Theorem~\ref{thm:temporal-evidence} and the pre-registration \cite{QG190}, the
prediction values were fixed in advance. Hence their confirmation is temporally
independent evidence, not post-hoc fitting.
\end{proof}

\section{Failure Conditions}
\label{sec:failure-conditions}

\begin{theorem}[The failure conditions]
\label{thm:failure-conditions}
The theory specifies explicit failure conditions: observations that would force a revision.
\end{theorem}

\begin{proof}
The prediction registry \cite{QG190} records the falsification condition of each frozen
prediction: for P1, no signal in statistically sensitive searches of the $99$--$114$ GeV
window; for P2, a measured upper limit below $2.02$ meV; for P3, a sensitive search
excluding any frozen rung. These conditions, documented in the outcome dashboard
\cite{QG202}, specify exactly which observations would contradict the theory. Hence the
failure conditions are explicit.
\end{proof}

\begin{corollary}[Revisions are possible]
\label{cor:revisions-possible}
The theory is revisable: a confirmed contradiction of a pre-registered prediction would
force a revision of the corresponding sector of the theory.
\end{corollary}

\begin{proof}
By Theorem~\ref{thm:failure-conditions}, the failure conditions specify the contradicting
observations. A confirmed contradiction would force a revision, exactly as a scientific
theory must allow. Hence the theory is revisable.
\end{proof}

\section{Scope of Validity}
\label{sec:scope}

\begin{theorem}[What the theory claims]
\label{thm:claims}
The theory claims: the derivation of its physics observables from the canonical hierarchy;
the pre-registered predictions; and the classification of its boundaries and validation
items.
\end{theorem}

\begin{proof}
The derivation chapters establish the observables; the pre-registration \cite{QG190}
establishes the predictions; Chapter~\ref{chap:boundaries} establishes the boundary
classification. Hence the theory's claims are exactly these.
\end{proof}

\begin{theorem}[What the theory does not claim]
\label{thm:non-claims-scope}
The theory does not claim: an open derivation frontier; a derivation of the hosted
standard-model Lagrangian; a derivation of the exact Bekenstein coefficient; or any result
beyond the accepted derivations.
\end{theorem}

\begin{proof}
There is no open derivation frontier (Corollary~\ref{cor:no-open-derivation}); the SM
dynamics is hosted (Chapter~\ref{chap:boundaries}, Theorem~\ref{thm:sm-hosted}); the
Bekenstein coefficient requires the imported $2\pi$ factor (Chapter~\ref{chap:boundaries},
Theorem~\ref{thm:bekenstein-2pi}); and the monograph is assembly-only
\cite{MONO007}. Hence the non-claims are explicit.
\end{proof}

\begin{corollary}[The scope is exact]
\label{cor:scope-exact}
The scope of validity of the theory is exact: it claims its derivations and predictions,
and discloses its boundaries and validation items, with no excess.
\end{corollary}

\begin{proof}
By Theorem~\ref{thm:claims} and Theorem~\ref{thm:non-claims-scope}, the claims and
non-claims are explicit and exhaustive. Hence the scope of validity is exact.
\end{proof}

\section{Future Experimental Program}
\label{sec:future-program}

\begin{theorem}[The validation roadmap]
\label{thm:validation-roadmap}
The future experimental program of the theory is a validation roadmap: the confirmation of
the derived predictions by the scheduled experiments.
\end{theorem}

\begin{proof}
The roadmap follows the validation targets. The oblique parameters $S,T,U$ require only
re-analysis as the EW fit tightens (low priority). The electron g-2 $a_e$ requires the
lattice-limited measurement (medium priority). The Majorana character $0\nu\beta\beta$
requires the next-generation experiments nEXO and LEGEND-1000 (high priority), which reach
the $0.036$--$0.156$ eV window containing $m_{\beta\beta} = 2.02$ meV \cite{VALID001}. The
pre-registered P1 (106 GeV) awaits the HL-LHC, and P3 awaits further high-energy
confirmation \cite{QG199,QG200}. Hence the future program is a validation roadmap.
\end{proof}

\begin{corollary}[The roadmap is experimental only]
\label{cor:roadmap-experimental}
The future program contains no derivation tasks: every observable is derived, and the
roadmap schedules only experimental validation.
\end{corollary}

\begin{proof}
By Theorem~\ref{thm:validation-roadmap}, the roadmap schedules the experiments; by
Corollary~\ref{cor:no-open-derivation}, there are no derivation tasks. Hence the future
program is experimental only.
\end{proof}

\begin{remark}[Publication-safe wording]
Consistent with MONO007 \cite{MONO007} and MONO110A \cite{MONO110A}, this chapter makes the
frontier explicit and experimental: no open derivation frontier is claimed, the validation
is clearly separated from derivation, the falsification paths are documented, and no claim
exceeds the accepted results. The core monograph is thereby closed on a publication-safe
basis.
\end{remark}

\section{Conclusion}
\label{sec:frontier-conclusion}

\begin{theorem}[The core monograph is closed]
\label{thm:core-closed}
The core of the physics-focused monograph is closed: the physics derivation is complete
within the canonical hierarchy, the boundaries are disclosed, and the frontier is an
explicit, experimental validation program.
\end{theorem}

\begin{proof}
The physics derivation is complete (the derivation chapters); the boundaries are disclosed
(Chapter~\ref{chap:boundaries}); the frontier is experimental
(Corollary~\ref{cor:no-open-derivation}, Theorem~\ref{thm:validation-roadmap}). Hence the
core monograph is closed on a publication-safe basis.
\end{proof}

\begin{corollary}[The publication is ready]
\label{cor:publication-ready}
The physics-focused monograph of The Actualization Theory is publication-ready: it derives
the physics, discloses the boundaries, and documents the experimental frontier.
\end{corollary}

\begin{proof}
By Theorem~\ref{thm:core-closed}, the core is complete, bounded, and frontier-explicit;
by MONO110A \cite{MONO110A}, the structure is physics-focused and publication-ready. Hence
the publication is ready.
\end{proof}

\section{Summary}
\label{sec:frontier-summary}

This chapter has closed the core monograph. We have proved:

\begin{enumerate}
\item \textbf{Falsifiability} (Theorem~\ref{thm:falsifiability},
Corollary~\ref{cor:falsifiability-strength}): the theory exposes its failure modes and is
therefore scientific;
\item \textbf{Derivation vs validation} (Theorem~\ref{thm:derived-vs-validated},
Corollary~\ref{cor:no-open-derivation}): derivation and validation are distinct, and there
is no open derivation frontier;
\item \textbf{Validation targets} (Theorem~\ref{thm:precision-validation},
Theorem~\ref{thm:pre-registered}, Corollary~\ref{cor:targets-derived}): $S,T,U$, $a_e$,
$0\nu\beta\beta$, P1, P2, and P3 are derived predictions awaiting experiment;
\item \textbf{Independent temporal evidence} (Theorem~\ref{thm:temporal-evidence},
Corollary~\ref{cor:temporal-independence}): the pre-registered predictions provide
temporally independent evidence;
\item \textbf{Failure conditions} (Theorem~\ref{thm:failure-conditions},
Corollary~\ref{cor:revisions-possible}): the falsification conditions are explicit and the
theory is revisable;
\item \textbf{Scope of validity} (Theorem~\ref{thm:claims},
Theorem~\ref{thm:non-claims-scope}, Corollary~\ref{cor:scope-exact}): the theory claims its
derivations and predictions, and discloses its boundaries;
\item \textbf{Future program} (Theorem~\ref{thm:validation-roadmap},
Corollary~\ref{cor:roadmap-experimental}): the future program is a validation roadmap, with
no derivation tasks;
\item \textbf{Conclusion} (Theorem~\ref{thm:core-closed},
Corollary~\ref{cor:publication-ready}): the core monograph is closed and publication-ready.
\end{enumerate}

The frontier of The Actualization Theory is therefore explicit and experimental: the derived
predictions --- $S,T,U$, the electron g-2, the Majorana character, P1, P2, and the P3
ladder --- await their validation, with explicit falsification conditions, temporally
independent evidence, and a validation roadmap. The theory claims no open derivation
frontier; it discloses its boundaries; and it is falsifiable. The core of the
physics-focused monograph is closed and publication-ready.

\begin{thebibliography}{99}
\bibitem{QG190} TQM-QG Phase 190, \emph{Pre-Registered Predictions}, Docs/Research.
\bibitem{QG199} TQM-QG Phase 199, \emph{P1 Evidence Update}, Docs/Research.
\bibitem{QG200} TQM-QG Phase 200, \emph{Sector Ladder Evidence Audit}, Docs/Research.
\bibitem{QG201} TQM-QG Phase 201, \emph{Ladder Statistics Audit}, Docs/Research.
\bibitem{QG202} TQM-QG Phase 202, \emph{Prediction Outcome Dashboard}, Docs/Research.
\bibitem{QG203} TQM-QG Phase 203, \emph{Absolute Neutrino Mass Origin}, Docs/Research.
\bibitem{QG299} TQM-QG Phase 299, \emph{Remaining Frontier Re-audit}, Docs/Research.
\bibitem{VALID001} TQM-VALID001, \emph{Remaining Untested Items Audit}, Docs/Research.
\bibitem{MONO007} TQM-MONO007, \emph{Referee-Readiness Resolution}, Docs/Zenodo/Monograph.
\bibitem{MONO110A} TQM-MONO110A, \emph{Physics-Focused Monograph Re-audit}, Docs/Research.
\end{thebibliography}
 after the preamble
%         that defines the theorem environments (or include via the master file).
% ======================================================================

% ---- Theorem environments (define in the master preamble if not present) ----
% \newtheorem{definition}{Definition}[chapter]
% \newtheorem{lemma}[definition]{Lemma}
% \newtheorem{theorem}[definition]{Theorem}
% \newtheorem{corollary}[definition]{Corollary}
% \newtheorem{remark}[definition]{Remark}

\chapter{Frontier and Falsification}
\label{chap:frontier}

\section{Introduction}
\label{sec:frontier-introduction}

The preceding chapters derived the physics of The Actualization Theory and documented its
boundaries. This chapter closes the core monograph by addressing the \emph{frontier}: the
experimental validation of the derived predictions, the independent temporal evidence, the
explicit failure conditions, and the scope of validity. The chapter's purpose is to make the
theory \emph{publication-safe} --- to state clearly what would confirm it, what would
falsify it, and what it does and does not claim.

The central content is: the scientific-falsifiability requirement \cite{QG299}; the
distinction between derivation and validation; the current validation targets --- the
oblique parameters $S,T,U$, the electron g-2 $a_e$, the Majorana character $0\nu\beta\beta$
\cite{VALID001}, and the pre-registered predictions P1 (106 GeV), P2 ($0\nu\beta\beta$), and the P3
sector ladder \cite{QG190,QG199,QG200,QG201,QG202,QG203}; the independent temporal evidence;
the explicit failure conditions; the scope of validity; and the future experimental
program.

Following the physics-focused re-audit \cite{MONO110A}, this chapter claims \emph{no open
derivation frontier}: every observable is derived, and the frontier is experimental. The
experimental validation is clearly separated from the derivation, the falsification paths
are explicit, and the wording is referee-safe and publication-safe \cite{MONO007}. We use
only accepted canonical results and introduce no new physics and no speculation.

\section{Scientific Falsifiability}
\label{sec:falsifiability}

\begin{definition}[Falsifiability]
\label{def:falsifiability}
A scientific theory is \emph{falsifiable} if it specifies, in advance, observations that
would contradict it: the theory must expose its failure modes.
\end{definition}

\begin{theorem}[A scientific theory must expose its failure modes]
\label{thm:falsifiability}
The Actualization Theory is falsifiable: it specifies explicit, pre-registered
observations whose absence or contradiction would force a revision.
\end{theorem}

\begin{proof}
The post-frontier audit \cite{QG299} classifies the remaining items into open, partial,
boundary, and methodology, with the pre-registered predictions carrying explicit
falsification conditions. The prediction-outcome dashboard \cite{QG202} records, for each
prediction, its frozen value, its current evidence, and its falsification condition.
Hence the theory exposes its failure modes in advance: it is falsifiable.
\end{proof}

\begin{corollary}[Falsifiability is not a weakness]
\label{cor:falsifiability-strength}
The exposure of failure modes is a strength, not a weakness: it is the property that makes
the theory scientific and its claims testable.
\end{corollary}

\begin{proof}
By Theorem~\ref{thm:falsifiability}, the theory specifies its failure modes; by
Definition~\ref{def:falsifiability}, this is the defining property of a falsifiable theory.
Hence the explicit falsification paths are a scientific strength.
\end{proof}

\section{Experimental Validation Status}
\label{sec:validation-status}

\begin{definition}[Derivation vs validation]
\label{def:derivation-validation}
\emph{Derivation} is the derivation of an observable from the theory's principles; 
\emph{validation} is the confirmation of that derived value by experiment. The two are
distinct: a derived prediction may await validation without any derivation gap.
\end{definition}

\begin{theorem}[Derived and validated are distinct]
\label{thm:derived-vs-validated}
The status of an observable is the pair (derived, validated): a quantity may be derived
without being validated, and its derivation is complete before its validation is possible.
\end{theorem}

\begin{proof}
Following VALID001 \cite{VALID001}, the quantities $S,T,U$, the electron g-2 $a_e$, and the
Majorana character $0\nu\beta\beta$ are status A --- derived prediction exists, with no
missing derivation steps --- and belong to the experimental-validation category. Their
derivation is complete; their validation is the open experimental task. Hence derivation and
validation are distinct, and a derived-but-unvalidated quantity is not an open derivation
issue.
\end{proof}

\begin{corollary}[No open derivation frontier]
\label{cor:no-open-derivation}
The theory claims no open derivation frontier: every observable of the theory is derived,
and the frontier is experimental validation, not derivation.
\end{corollary}

\begin{proof}
By Theorem~\ref{thm:derived-vs-validated}, the derived quantities are complete, with the
validation tasks ahead. The reconstruction audit establishes that the major results are
reconstructed from the minimal hierarchy. Hence there is no open derivation frontier; the
frontier is experimental \cite{MONO110A}.
\end{proof}

\section{Current Validation Targets}
\label{sec:validation-targets}

\begin{theorem}[The precision predictions await validation]
\label{thm:precision-validation}
The oblique parameters $S,T,U$, the electron g-2 $a_e$, and the Majorana character
$0\nu\beta\beta$ are derived predictions awaiting experimental validation.
\end{theorem}

\begin{proof}
Following VALID001 \cite{VALID001}: $S = 0.0421$, $T = 0.0842$, $U = 0$ are derived and
consistent with the EW global fit; $a_e = 1.159655\times10^{-3}$ matches experiment within
$0.0003\%$; $m_{\beta\beta} = 2.02\times10^{-3}$ eV is within limits and in reach of
next-generation experiments. All three are status A, with their experimental validation the
open task. Hence they are current validation targets.
\end{proof}

\begin{theorem}[The pre-registered predictions]
\label{thm:pre-registered}
The theory carries three pre-registered predictions with frozen values and explicit
falsification conditions:
\begin{enumerate}
\item \textbf{P1} --- a 106 GeV resonance: central value $106.39$ GeV, window
$99$--$114$ GeV;
\item \textbf{P2} --- the $0\nu\beta\beta$ effective Majorana mass $m_{\beta\beta} =
2.02$ meV;
\item \textbf{P3} --- the sector-ladder spectrum: nine resonances
$106.39 \to 263.43$ GeV with a $15.20$ GeV width scale.
\end{enumerate}
\end{theorem}

\begin{proof}
The predictions are pre-registered \cite{QG190}, with their values frozen in the immutable
prediction registry and their evidence tracked in the outcome dashboard \cite{QG202}. P1 is
pending (window $99$--$114$ GeV neither confirmed nor excluded) \cite{QG199}; P2 is pending
(below current sensitivity, awaiting nEXO/LEGEND-1000); P3 is supported (the $151.98$ GeV
rung matches a $\sim152$ GeV excess at moderate significance) \cite{QG200,QG201}. The
absolute-mass law \cite{QG203} anchors the sector structure.
\end{proof}

\begin{corollary}[The targets are derived, not open issues]
\label{cor:targets-derived}
The validation targets are derived predictions awaiting experiment; they are not open
derivation issues and not boundaries.
\end{corollary}

\begin{proof}
By Theorem~\ref{thm:precision-validation} and Theorem~\ref{thm:pre-registered}, all
targets are derived predictions with frozen values; by Chapter~\ref{chap:boundaries}
(Corollary~\ref{cor:no-open-validation}), validation items are distinct from open
derivation issues. Hence the targets are validation items.
\end{proof}

\section{Independent Temporal Evidence}
\label{sec:temporal-evidence}

\begin{definition}[Independent temporal evidence]
\label{def:temporal-evidence}
\emph{Independent temporal evidence} is the empirical program of the theory that does not
depend on the already-derived observables: the confirmation of the pre-registered
predictions by observations made after the predictions were frozen.
\end{definition}

\begin{theorem}[The independent evidence program is the remaining external task]
\label{thm:temporal-evidence}
The independent temporal evidence program is the remaining external empirical task of the
theory: the confirmation of the pre-registered predictions by new, temporally-ordered
observations.
\end{theorem}

\begin{proof}
The pre-registered predictions \cite{QG190} were frozen before their observation; the
outcome dashboard \cite{QG202} tracks their evidence. The confirmations (or
contradictions) of P1, P2, and P3 by new observations constitute the independent temporal
evidence --- evidence obtained after the predictions were fixed, hence independent of the
derivation. This is the remaining external empirical program \cite{QG299}.
\end{proof}

\begin{corollary}[Temporal independence is preserved]
\label{cor:temporal-independence}
The pre-registered predictions preserve temporal independence: their values were frozen
before the observations, so their confirmation is genuine independent evidence.
\end{corollary}

\begin{proof}
By Theorem~\ref{thm:temporal-evidence} and the pre-registration \cite{QG190}, the
prediction values were fixed in advance. Hence their confirmation is temporally
independent evidence, not post-hoc fitting.
\end{proof}

\section{Failure Conditions}
\label{sec:failure-conditions}

\begin{theorem}[The failure conditions]
\label{thm:failure-conditions}
The theory specifies explicit failure conditions: observations that would force a revision.
\end{theorem}

\begin{proof}
The prediction registry \cite{QG190} records the falsification condition of each frozen
prediction: for P1, no signal in statistically sensitive searches of the $99$--$114$ GeV
window; for P2, a measured upper limit below $2.02$ meV; for P3, a sensitive search
excluding any frozen rung. These conditions, documented in the outcome dashboard
\cite{QG202}, specify exactly which observations would contradict the theory. Hence the
failure conditions are explicit.
\end{proof}

\begin{corollary}[Revisions are possible]
\label{cor:revisions-possible}
The theory is revisable: a confirmed contradiction of a pre-registered prediction would
force a revision of the corresponding sector of the theory.
\end{corollary}

\begin{proof}
By Theorem~\ref{thm:failure-conditions}, the failure conditions specify the contradicting
observations. A confirmed contradiction would force a revision, exactly as a scientific
theory must allow. Hence the theory is revisable.
\end{proof}

\section{Scope of Validity}
\label{sec:scope}

\begin{theorem}[What the theory claims]
\label{thm:claims}
The theory claims: the derivation of its physics observables from the canonical hierarchy;
the pre-registered predictions; and the classification of its boundaries and validation
items.
\end{theorem}

\begin{proof}
The derivation chapters establish the observables; the pre-registration \cite{QG190}
establishes the predictions; Chapter~\ref{chap:boundaries} establishes the boundary
classification. Hence the theory's claims are exactly these.
\end{proof}

\begin{theorem}[What the theory does not claim]
\label{thm:non-claims-scope}
The theory does not claim: an open derivation frontier; a derivation of the hosted
standard-model Lagrangian; a derivation of the exact Bekenstein coefficient; or any result
beyond the accepted derivations.
\end{theorem}

\begin{proof}
There is no open derivation frontier (Corollary~\ref{cor:no-open-derivation}); the SM
dynamics is hosted (Chapter~\ref{chap:boundaries}, Theorem~\ref{thm:sm-hosted}); the
Bekenstein coefficient requires the imported $2\pi$ factor (Chapter~\ref{chap:boundaries},
Theorem~\ref{thm:bekenstein-2pi}); and the monograph is assembly-only
\cite{MONO007}. Hence the non-claims are explicit.
\end{proof}

\begin{corollary}[The scope is exact]
\label{cor:scope-exact}
The scope of validity of the theory is exact: it claims its derivations and predictions,
and discloses its boundaries and validation items, with no excess.
\end{corollary}

\begin{proof}
By Theorem~\ref{thm:claims} and Theorem~\ref{thm:non-claims-scope}, the claims and
non-claims are explicit and exhaustive. Hence the scope of validity is exact.
\end{proof}

\section{Future Experimental Program}
\label{sec:future-program}

\begin{theorem}[The validation roadmap]
\label{thm:validation-roadmap}
The future experimental program of the theory is a validation roadmap: the confirmation of
the derived predictions by the scheduled experiments.
\end{theorem}

\begin{proof}
The roadmap follows the validation targets. The oblique parameters $S,T,U$ require only
re-analysis as the EW fit tightens (low priority). The electron g-2 $a_e$ requires the
lattice-limited measurement (medium priority). The Majorana character $0\nu\beta\beta$
requires the next-generation experiments nEXO and LEGEND-1000 (high priority), which reach
the $0.036$--$0.156$ eV window containing $m_{\beta\beta} = 2.02$ meV \cite{VALID001}. The
pre-registered P1 (106 GeV) awaits the HL-LHC, and P3 awaits further high-energy
confirmation \cite{QG199,QG200}. Hence the future program is a validation roadmap.
\end{proof}

\begin{corollary}[The roadmap is experimental only]
\label{cor:roadmap-experimental}
The future program contains no derivation tasks: every observable is derived, and the
roadmap schedules only experimental validation.
\end{corollary}

\begin{proof}
By Theorem~\ref{thm:validation-roadmap}, the roadmap schedules the experiments; by
Corollary~\ref{cor:no-open-derivation}, there are no derivation tasks. Hence the future
program is experimental only.
\end{proof}

\begin{remark}[Publication-safe wording]
Consistent with MONO007 \cite{MONO007} and MONO110A \cite{MONO110A}, this chapter makes the
frontier explicit and experimental: no open derivation frontier is claimed, the validation
is clearly separated from derivation, the falsification paths are documented, and no claim
exceeds the accepted results. The core monograph is thereby closed on a publication-safe
basis.
\end{remark}

\section{Conclusion}
\label{sec:frontier-conclusion}

\begin{theorem}[The core monograph is closed]
\label{thm:core-closed}
The core of the physics-focused monograph is closed: the physics derivation is complete
within the canonical hierarchy, the boundaries are disclosed, and the frontier is an
explicit, experimental validation program.
\end{theorem}

\begin{proof}
The physics derivation is complete (the derivation chapters); the boundaries are disclosed
(Chapter~\ref{chap:boundaries}); the frontier is experimental
(Corollary~\ref{cor:no-open-derivation}, Theorem~\ref{thm:validation-roadmap}). Hence the
core monograph is closed on a publication-safe basis.
\end{proof}

\begin{corollary}[The publication is ready]
\label{cor:publication-ready}
The physics-focused monograph of The Actualization Theory is publication-ready: it derives
the physics, discloses the boundaries, and documents the experimental frontier.
\end{corollary}

\begin{proof}
By Theorem~\ref{thm:core-closed}, the core is complete, bounded, and frontier-explicit;
by MONO110A \cite{MONO110A}, the structure is physics-focused and publication-ready. Hence
the publication is ready.
\end{proof}

\section{Summary}
\label{sec:frontier-summary}

This chapter has closed the core monograph. We have proved:

\begin{enumerate}
\item \textbf{Falsifiability} (Theorem~\ref{thm:falsifiability},
Corollary~\ref{cor:falsifiability-strength}): the theory exposes its failure modes and is
therefore scientific;
\item \textbf{Derivation vs validation} (Theorem~\ref{thm:derived-vs-validated},
Corollary~\ref{cor:no-open-derivation}): derivation and validation are distinct, and there
is no open derivation frontier;
\item \textbf{Validation targets} (Theorem~\ref{thm:precision-validation},
Theorem~\ref{thm:pre-registered}, Corollary~\ref{cor:targets-derived}): $S,T,U$, $a_e$,
$0\nu\beta\beta$, P1, P2, and P3 are derived predictions awaiting experiment;
\item \textbf{Independent temporal evidence} (Theorem~\ref{thm:temporal-evidence},
Corollary~\ref{cor:temporal-independence}): the pre-registered predictions provide
temporally independent evidence;
\item \textbf{Failure conditions} (Theorem~\ref{thm:failure-conditions},
Corollary~\ref{cor:revisions-possible}): the falsification conditions are explicit and the
theory is revisable;
\item \textbf{Scope of validity} (Theorem~\ref{thm:claims},
Theorem~\ref{thm:non-claims-scope}, Corollary~\ref{cor:scope-exact}): the theory claims its
derivations and predictions, and discloses its boundaries;
\item \textbf{Future program} (Theorem~\ref{thm:validation-roadmap},
Corollary~\ref{cor:roadmap-experimental}): the future program is a validation roadmap, with
no derivation tasks;
\item \textbf{Conclusion} (Theorem~\ref{thm:core-closed},
Corollary~\ref{cor:publication-ready}): the core monograph is closed and publication-ready.
\end{enumerate}

The frontier of The Actualization Theory is therefore explicit and experimental: the derived
predictions --- $S,T,U$, the electron g-2, the Majorana character, P1, P2, and the P3
ladder --- await their validation, with explicit falsification conditions, temporally
independent evidence, and a validation roadmap. The theory claims no open derivation
frontier; it discloses its boundaries; and it is falsifiable. The core of the
physics-focused monograph is closed and publication-ready.

\begin{thebibliography}{99}
\bibitem{QG190} TQM-QG Phase 190, \emph{Pre-Registered Predictions}, Docs/Research.
\bibitem{QG199} TQM-QG Phase 199, \emph{P1 Evidence Update}, Docs/Research.
\bibitem{QG200} TQM-QG Phase 200, \emph{Sector Ladder Evidence Audit}, Docs/Research.
\bibitem{QG201} TQM-QG Phase 201, \emph{Ladder Statistics Audit}, Docs/Research.
\bibitem{QG202} TQM-QG Phase 202, \emph{Prediction Outcome Dashboard}, Docs/Research.
\bibitem{QG203} TQM-QG Phase 203, \emph{Absolute Neutrino Mass Origin}, Docs/Research.
\bibitem{QG299} TQM-QG Phase 299, \emph{Remaining Frontier Re-audit}, Docs/Research.
\bibitem{VALID001} TQM-VALID001, \emph{Remaining Untested Items Audit}, Docs/Research.
\bibitem{MONO007} TQM-MONO007, \emph{Referee-Readiness Resolution}, Docs/Zenodo/Monograph.
\bibitem{MONO110A} TQM-MONO110A, \emph{Physics-Focused Monograph Re-audit}, Docs/Research.
\end{thebibliography}
